% TA-PFC-E1 — Luis Yépez. Identificación de perfiles.
% Fragmento para insertar después de \subsection{Identificación de perfiles de usuarios y actores}.

Para el primer avance se identifican cuatro perfiles de usuarios previstos. Sus responsabilidades se proponen a partir de la visión actualizada del proyecto y deberán contrastarse con la institución y con las actividades de elicitación; no representan permisos concedidos ni entrevistas realizadas.

\begin{description}
  \item[Estudiante.] Persona que recibe el acompañamiento académico. Se prevé que consulte la disponibilidad de tutorías, solicite sesiones, confirme su participación y revise su propio historial tutorial. Podría visualizar indicadores de su evolución académica y recomendaciones de tutorías cuando estén habilitados para su perfil, sin acceder a los datos de otros estudiantes.

  \item[Docente tutor.] Responsable de atender tutorías y registrar las actividades de acompañamiento. Se propone que gestione horarios o solicitudes, documente sesiones y asistencias y consulte, únicamente para los estudiantes a su cargo y según autorización, los indicadores necesarios para orientar acciones tutoriales. También participará en la validación de los criterios de alerta y en la revisión de recomendaciones.

  \item[Administrador del sistema.] Responsable de la administración técnica de la aplicación. Se propone que gestione cuentas, roles, permisos y configuración operativa, y atienda incidencias de acceso. Este perfil no debería modificar calificaciones ni decidir por sí mismo qué estudiantes requieren una intervención académica.

  \item[Coordinador académico.] Responsable previsto para supervisar la ejecución del proceso tutorial. Se propone que consulte reportes de seguimiento e indicadores pertinentes, revise casos que requieren atención y coordine intervenciones con los docentes tutores, dentro de las atribuciones que la institución confirme. La autoridad para aprobar criterios de alerta y requisitos aún debe verificarse.
\end{description}

Estos perfiles son \emph{usuarios previstos} y, a la vez, partes interesadas en las necesidades que aportan. Otros actores institucionales podrían ser relevantes aunque no utilicen la aplicación directamente, por ejemplo quienes administren la fuente autorizada de calificaciones o las políticas de tratamiento de datos personales. Su existencia, participación y nivel de decisión deben confirmarse, sin asumir integraciones institucionales no autorizadas.
